\documentclass[11pt]{article}
\newcommand{\runningtitle}{Prime sums and individual Hecke factors}
\usepackage{lmodern}
\pdfmapfile{+lm.map}
\usepackage[a4paper,top=24mm,bottom=25mm,left=28mm,right=28mm]{geometry}
\usepackage{amsmath,amssymb,booktabs,array,mathtools,enumitem,xcolor,fancyhdr,hyperref}
\definecolor{ink}{HTML}{17364B}
\definecolor{teal}{HTML}{00685E}
\hypersetup{colorlinks=true,linkcolor=teal,urlcolor=teal,citecolor=teal,pdfauthor={OpenAI}}
\urlstyle{same}
\setlength{\parskip}{3.1pt plus 1pt}
\setlength{\parindent}{0pt}
\setlist{itemsep=1pt,topsep=3pt}
\setlength{\headheight}{15pt}
\pagestyle{fancy}\fancyhf{}
\fancyhead[L]{\small\sffamily\color{ink}\runningtitle}
\fancyhead[R]{\small\sffamily\color{ink}Research note}
\fancyfoot[R]{\small\thepage}
\fancyfoot[L]{\small\sffamily\color{ink}Transfers from the Dirichlet GRH proof}
\renewcommand{\headrulewidth}{0.3pt}
\newcommand{\headline}[3]{%
 {\small\sffamily\bfseries\color{teal}#1\par}
 \vspace{6pt}{\raggedright\LARGE\bfseries\color{ink}#2\par}
 \vspace{7pt}{\small\sffamily #3\par}\vspace{8pt}
 {\color{teal}\hrule height 0.9pt}\vspace{8pt}}
\newcommand{\grhlink}[2]{\href{https://github.com/lsernau/grh-lean/blob/72fead102f663d5c7a407964a13dc0cf9d977207/grh/#1}{#2}}
\newcommand{\incominglink}[2]{\href{https://github.com/VladimirReshetnikov/ProveIt/blob/274841c1cb137aafbb60ab5f519c51fd42eee98e/docs/incoming/#1}{#2}}
\setcounter{tocdepth}{1}
\setlength{\emergencystretch}{2em}
\raggedbottom

\begin{document}
\headline{ANALYTIC NUMBER THEORY}{Prime sums that isolate a Hecke factor}{One-sided criteria beyond Dirichlet characters and an excursion bound \hfill 25 November 2026}
A one-sided bound on a prime sum can locate zeros only if the relevant
zero contributes a pole with a controlled sign. In the readable Lean
\grhlink{src/New/Reduction/Zeros/PrimeError.lean}{prime-error computation}
the residues depend on the difference of two zero multiplicities.
Its \grhlink{src/New/Reduction/Zeros/Criterion.lean}{criterion}
chooses a quadratic twist that has no competing zero in a prescribed
rectangle. We instead choose weights on ideal classes so that unwanted
unknown factors disappear from the logarithmic derivative. The remaining
residues have a common sign. The result identifies one prime estimate
equivalent to RH for a selected Hecke function, including functions of
genuinely nonabelian origin over $\mathbb Q$.

For example, over $\mathbb Q(\sqrt{-95})$, three positive definite
binary quadratic forms with weights $1/2,\sqrt{2}/2,1$ suffice to test
one faithful order-eight character pair. The other faithful pair is
eliminated identically. Rational-valued weights cannot achieve this
specific separation. A second example uses a single quadratic form at
each step in the Gaussian ring class tower.

The \incominglink{grh_sparse_singular_moduli.zip}{prior sparse singular
moduli report} already gives the aggregate commutator test and a
dihedral order-eight field example. Our contributions are the
factor-specific optimization, its precise rationality obstruction and
the explicit tower formulation. These are equivalences for open zero
problems. The estimates required for a new Hecke or Dedekind RH remain
unproved here. The last section obtains a numerical-width bound for the
lower limit of a different observable, the repository's signed
modulus-one remainder; it uses RH and Dirichlet GRH.

\section{One-sided tests for class-group Hecke functions}
Let $K$ be an imaginary quadratic field, $C$ its ideal class group,
$h=|C|$, and $B\le C$. The characters trivial on $B$ form $B^\perp$.
Suppose that their Hecke functions $L_K(s,\eta)$ have no zeros in
$\Re s>1/2$. We call these the known factors. For $B=C^2$, all such
characters have order at most two. Kronecker's genus factorization expresses
their $L$-functions as products of two Dirichlet $L$-functions, including
$\zeta_K=\zeta L(\cdot,\chi_K)$ for the trivial character. Thus the
Dirichlet theorem in the repository supplies precisely this hypothesis.
Hecke's continuation, nonvanishing at $1$ for nontrivial characters and the
functional equation apply to the other characters as well.

For a real inversion-invariant weight $w:C\to[0,1]$, put
\[
 a_\chi={1\over h}\sum_{c\in C} w(c)\overline{\chi(c)},\qquad
 \Theta_w(x)=\sum_{\substack{p\le x\\p\ \mathrm{split\ in}\ K}}
     w([\mathfrak p])\log p,\qquad \delta=a_1/2 .
\]
Either prime ideal above $p$ gives the same weight, and each rational prime
is counted once. The coefficients $a_\chi$ are real.
Let $S$ be a nonempty inversion-invariant subset of
$\widehat C\setminus B^\perp$. Require $a_\chi<0$ on $S$ and
$a_\chi=0$ off $S\cup B^\perp$.

\paragraph{Criterion.}
The following three statements are equivalent, with fixed field and weight.
\begin{enumerate}
\item Each $L_K(s,\chi)$ with $\chi\in S$ satisfies RH.
\item For each $0<\epsilon<1/2$ and all sufficiently large $x$,
$\Theta_w(x)-\delta x\le M_\epsilon x^{1/2+\epsilon}$ for some $M_\epsilon$.
\item For every such $\epsilon$ and all sufficiently large $x$,
\[
 R_w(x)=\sum_{\substack{p\le x\\p\ \mathrm{split\ in}\ K}}
 w([\mathfrak p])\log p\,(p^{-1}-x^{-1})
       -\delta\log x-c_w \ \ge -M_\epsilon x^{-1/2+\epsilon}.
\]
\end{enumerate}
Here $P_w(s)=\sum_{p\ \mathrm{split}}w([\mathfrak p])\log(p)p^{-s}$ and
$c_w=\lim_{s\to1}(P_w(s)-\delta/(s-1))-\delta$.

\paragraph{Proof.}
Comparing Euler factors for $\Re s>1$ gives
\[
 P_w(s)=-\frac12\sum_\chi a_\chi\frac{L'_K}{L_K}(s,\chi)+H(s)
\]
with $H$ holomorphic on $\Re s>1/2$.
The two degree-one primes above a split prime contribute equally in the
weighted sum. Inert primes have norm $p^2$, the terms for powers of exponent
at least two converge absolutely and locally uniformly in that half-plane,
and ramified primes contribute finitely many locally holomorphic series.
The pole at $1$ has residue $\delta$. At a zero $\rho$, $1/2<\Re\rho<1$,
of the tested functions the residue of $P_w$ is exactly
\[
 d_\rho=\frac12\sum_{\chi\in S}(-a_\chi)\operatorname{ord}_\rho L_K(s,\chi)>0.
\]
An untested function vanishing at the same point does not affect this
formula. In particular no assertion about independence or simplicity of
zeros is needed.

Write $E(x)=\Theta_w(x)-\delta x$ and use the transform
$\mathcal M F(s)=\int_1^\infty F(x)x^{-s-1}\,dx$. Termwise integration
in $\Re s>1$ gives
\begin{align*}
 \mathcal M E(s)&=P_w(s)/s-\delta/(s-1),\\
 \mathcal M(xR_w)(s)&=P_w(s)/(s(s-1))-\delta/(s-1)^2-c_w/(s-1).
\end{align*}
Both expressions are regular at $1$. Take $a=1/2+\epsilon$. Under the
second assertion, $M_\epsilon x^a-E(x)$ is eventually nonnegative.
Under the third, $M_\epsilon x^a+xR_w(x)$ is eventually nonnegative.
In each case the transform is meromorphic in $\Re s>a$, and every pole
on the real axis there would have a strictly negative residue, respectively
$-d_\rho/\rho$ and $d_\rho/(\rho(\rho-1))$.
Landau's theorem for the Mellin integral of an eventually nonnegative
function forces a singularity at any finite real abscissa of convergence.
This follows by expanding at a real point just to its right: positivity
of the moments and monotone convergence would extend the integral leftward
if the real abscissa were regular. A modification on a bounded interval
makes the integrand nonnegative without changing the poles.
A pole of negative residue at the abscissa is impossible since the
nonnegative integral would tend to minus infinity from the right. Thus the
abscissa is at most $a$ and the integral is holomorphic on $\Re s>a$.
Either proposed bound rules out all tested zeros there. Vary $a$ and use the
functional equation, with $S$ closed under inversion.

Conversely, the fixed-character prime ideal explicit formula under RH gives
\[\sum_{N\mathfrak p^k\le x}\chi(\mathfrak p)^k\log N\mathfrak p
=\mathbf1_{\chi=1}x+O_{K,\chi}(\sqrt{x}\log^2(2x))\] for each factor
with $a_\chi\ne0$. In the truncated explicit formula one takes height of
order $x$; zeros on the critical line contribute at most $\sqrt{x}$ times
$O_{K,\chi}(\log^2(2x))$ by the standard unit-height zero count.
Removing prime powers, inert primes and ramified primes costs
$O_K(\sqrt{x}\log^2(2x))$. Fourier inversion gives this same two-sided
error for $\Theta_w$. Finally the smoothed prime sum equals
$\int_1^x\Theta_w(t)t^{-2}\,dt$. The two-sided error makes
$\int_1^\infty E(t)t^{-2}\,dt$ converge and bounds its tail by
$O_{K,w}(x^{-1/2}\log^2(2x))$. Evaluation of $\mathcal ME$ at $1$
identifies the integral with $c_w$. This gives the third assertion.
\hfill$\square$

\paragraph{Principal genus.}
For $w=\mathbf1_{B\setminus\{1\}}$, orthogonality gives
$a_\eta=(|B|-1)/h$ for $\eta\in B^\perp$ and $a_\chi=-1/h$ otherwise.
The test therefore detects every remaining factor; its prime density is
$(|B|-1)/(2h)$. In a ring class field the analogous event is exactly
splitting completely in the fixed field of $B$ and failing to split
completely in the full field. This aggregate case specializes the earlier
ProveIt report's normal-subgroup criterion. Selecting a single pair of
factors requires a different weight.

\section{Isolating one dihedral factor}
Fix a character $\chi$ outside $B^\perp$ of order $n>2$.
Over an imaginary quadratic field $L_K(s,\chi)=L_K(s,\chi^{-1})$:
conjugation of ideals preserves norms and inverts their classes.
We seek a single bounded weight with $a_\chi=a_{\chi^{-1}}=-t<0$
and zero coefficients on every other unknown character. The preceding
criterion then tests just this Hecke function even if other untested
functions have off-line zeros.

\paragraph{Sharp coefficient and support.}
Put $v(c)=\Re\chi(c)$. For each coset $Q$ of $B$ let
$M_Q=\max_{c\in Q} v(c)$, $m_Q=\min_{c\in Q}v(c)$ and
$D=\max_Q(M_Q-m_Q)$. Then $D>0$. The largest possible $t$ is
$1/(2D)$, achieved by
\[
 w_\chi(c)=\frac{M_{cB}-v(c)}D .
\]
This choice also attains the smallest support by inclusion.
Indeed Fourier inversion gives $w(c)=u(cB)-2t v(c)$, where $u$ is a real,
inversion-invariant function on the quotient. The bounds $0\le w\le1$
hold precisely when $2tM_Q\le u(Q)\le1+2t m_Q$ for every $Q$.
Since $\chi|_B$ is nontrivial, $v$ is nonconstant on $B$. Taking
$u(Q)=2tM_Q$ makes the weight vanish at every coset maximum. All other
points have strictly positive weight for every feasible $u$ and every
$t>0$. The asserted sharpness follows.

\paragraph{Closed form for the principal genus.}
For $B=C^2$ the values of $t_{\max}$ are
\[
\begin{array}{c|ccc}
 n & n\ \text{odd} & 4\mid n & n\equiv2\pmod4\\ \hline
 t_{\max}&[2(1+\cos(\pi/n))]^{-1}&1/4&
                 [2(1+\cos(2\pi/n))]^{-1}.
\end{array}
\]
For odd $n$, the image of $B$ is all $n$th roots, with extreme real parts
$1,-\cos(\pi/n)$. For even $n$ each coset maps onto the even or odd
powers of a primitive $n$th root. The even powers include $\pm1$ when
$4\mid n$. For $n\equiv2\pmod4$ they form an odd regular polygon, and
the odd powers are its negative, giving the last entry.

\paragraph{Minimum coefficient field.}
For an exact test of one inverse pair, the smallest possible real number
field containing all weights is $F_n=\mathbb Q(\zeta_n+\zeta_n^{-1})$.
Indeed the optimal weight above already lies in $F_n$. Conversely let
$E$ be any number field containing all weights. Every automorphism of
$\overline{\mathbb Q}$ fixing $E$ induces a power map on $\mu_n$.
It permutes the nonzero Fourier coefficients belonging to primitive
powers of $\chi$. All these powers are nontrivial on $B$, since raising
a root of unity to a power coprime to $n$ preserves its order. Hence the
image of $\chi$ must be $\chi$ or $\chi^{-1}$. Consequently every
such automorphism fixes $\zeta_n+\zeta_n^{-1}$, so $F_n\subseteq E$.

\paragraph{When can the counted primes have weight one?}
For $B=C^2$, a rational-valued weight isolating this single pair exists
exactly for $n=3,4,6$. In all three cases a zero-one weight suffices.
To see necessity, an automorphism taking $\zeta_n$ to $\zeta_n^j$ sends
$a_\chi$ to $a_{\chi^j}$ for every $j$ coprime to $n$. A rational weight
therefore has nonzero coefficients throughout that orbit. Every member
still has order $n>2$ and thus lies outside the known characters. An
isolated pair forces $\varphi(n)\le2$, giving the listed orders. For
sufficiency write $\chi(c)=\exp(2\pi i j/n)$. Take the indicators of the
following exponents:
\[
 n=3:\ \{1,2\},\qquad n=4:\ \{2\},\qquad
 n=6:\ \{2,3,4\}.
\]
These are the weights $w_\chi$ above, as direct evaluation shows.
The classification concerns exact vanishing of the extraneous Fourier
coefficients. It allows no inference about unknown coincidences of zeros.

\paragraph{A pair that requires irrational weights.}
For $C=\mathbb Z/8\mathbb Z$ and $\chi(j)=\exp(2\pi i j/8)$, in order
$j=0,\ldots,7$ the optimal weight is
\[
 (0,0,1/2,\sqrt2/2,1,\sqrt2/2,1/2,0).
\]
Its negative coefficient on each of $\chi,\chi^{-1}$ equals $-1/4$,
and its rational-prime mean is $(2+\sqrt2)/16$. The primitive character
$\chi^3$ gives a second pair with no contribution to this test. Replacing
$j$ by $3j$ gives its own test. A rational weight could detect both pairs
together, but cannot distinguish them with this exact Fourier criterion.

A maximal-order realization is $K=\mathbb Q(\sqrt{-95})$.
The eight reduced primitive forms of discriminant $-95$ are
\[
 [1,1,24],\quad[2,\pm1,12],\quad[3,\pm1,8],\quad
 [4,\pm1,6],\quad[5,5,6].
\]
The list follows from $|b|\le a\le c$ and $a\le\sqrt{95/3}$, taking
$b\ge0$ when a boundary equality holds. Let $g=[2,1,12]$. The root
$b=1$ of $b^2\equiv-95$ modulo $4\cdot 2^k$ for $k=1,2,3$ gives
$g^2=[4,1,6]$ and $g^3=[8,1,3]=[3,-1,8]$ by ordinary form composition.
The square $[4,1,6]$ is distinct from its reduced inverse $[4,-1,6]$,
so $g$ has order eight. The only nonprincipal ambiguous class
$[5,5,6]$ is $g^4$. Thus the three distinct rational-prime representation
tests required by the displayed weight have forms and coefficients
\[
 \begin{array}{c|ccc}
 Q(x,y)&4x^2+xy+6y^2&3x^2+xy+8y^2&5x^2+5xy+6y^2\\\hline
 w&1/2&\sqrt2/2&1 .
 \end{array}
\]
Omit $p\mid95$. For prime inputs the three sets are disjoint: a split
prime determines exactly one unordered pair of inverse classes. Replacing
the middle form by $2x^2+xy+12y^2$ switches to the other faithful pair.

\section{One quadratic form for each layer above the Gaussian field}
\label{sec:gaussian}

There is a particularly simple application where a single
unweighted quadratic form tests a whole new layer. Put $K=\mathbb Q(i)$.
For $n\ge1$, let $K_n$ be its ring class field of conductor $2^n$ and set
\[
 Q_n(X,Y)=(2X+Y)^2+2^{2n-2}Y^2\qquad(n\ge2).
\]
For odd rational primes define
$\vartheta_n(x)=\sum_{p\le x,\ p=Q_n(a,b)}\log p$,
counting each represented prime once.

\paragraph{Tower criterion.}
Assume RH for $\zeta_{K_{n-1}}$. Then RH for $\zeta_{K_n}$ is
equivalent to the upper bound
\[
              \vartheta_n(x)\le 2^{-n}x+
                         O_{n,\epsilon}(x^{1/2+\epsilon})
              \quad\text{for every }\epsilon>0 .
\]
In particular these tests for all $n\ge3$ together are equivalent, under
Dirichlet GRH, to RH for the whole tower. Constants may depend on $n$.
The criterion at $n=3$ tests $4X^2+4XY+17Y^2$ with density $1/8$.
Here $K_3=\mathbb Q(i,\sqrt[4]{2})$ (Cox, \S9, ring class field
of $\mathbb Z[\sqrt{-64}]$), and $K_2=\mathbb Q(\zeta_8)$.
Thus the first test counts exactly the quartic nonresidues among
$p\equiv1\pmod8$ considered in the incoming sparse singular-moduli
report, with $2$ as the numerator of the quartic symbol.
The formula $Q_n$ extends this example to every subsequent layer.

\paragraph{Which primes are counted.}
The ring class number formula, including the factor two from the
Gaussian units, gives $[K_n:K]=2^{n-1}$.
The ring class reciprocity law says that an odd prime splits completely
in $K_n$ exactly when $p=a^2+2^{2n}b^2$ for integers $a,b$
(see Cox, \emph{Primes of the form $x^2+ny^2$}, \S\S7,9).
These facts also follow from the exact sequence with quotient
$(\mathbb Z[i]/2^n)^\times/((\mathbb Z/2^n)^\times\mathbb Z[i]^\times)$,
since $\mathbb Z[i]$ has class number one. The fields are nested by
the quotient maps of these groups. For $p\equiv1\pmod 4$ write
$p=a^2+c^2$ with $a$ odd and $c$ even. The pair of absolute values in
this order is unique by unique factorization in $\mathbb Z[i]$.
Thus complete splitting in $K_{n-1}$ and failure of complete splitting
in $K_n$ means $v_2(c)=n-1$. Equivalently $c=2^{n-1}Y$ with $Y$ odd,
and $a=2X+Y$ for some integer $X$. Conversely an odd prime of the displayed
form forces $Y$ odd and has this exact valuation. This proves the promised
description. The quadratic extension $K_2/K$ is Galois over $\mathbb Q$,
and its degree-four Galois group is abelian. Hence Dirichlet GRH supplies
its zeta RH by abelian factorization, starting the induction.

\paragraph{Why an upper bound suffices.}
For any fixed Galois field $M/\mathbb Q$, the series
$P_M(s)=\sum_{p\,\mathrm{split\ completely}}(\log p)p^{-s}$
equals $-[M:\mathbb Q]^{-1}\zeta_M'(s)/\zeta_M(s)$ up to a function
holomorphic on $\Re s>1/2$. Indeed the unramified degree-one prime
ideals occur there in full packets of $[M:\mathbb Q]$; all remaining
unramified terms have exponent at least two. The ramified Euler factors
are harmless in this region. The tested series is $P_{K_{n-1}}-P_{K_n}$
up to any finitely omitted primes. Its pole at 1 has residue
$2^{1-n}-2^{-n}=2^{-n}$. Under RH for the lower field, a zero of
the upper zeta function at $\beta+i\gamma$, $1/2<\beta<1$, contributes
a strictly positive residue equal to its multiplicity divided by $2^n$.
The unsmoothed Landau argument of the preceding sections gives the forward
implication from the bound to zero-freeness. Conversely RH for both
fields and the fixed-field prime ideal explicit formula give the
two-sided $O_n(\sqrt x\log^2(2x))$ error. This proof requires no
holomorphy conjecture for individual Artin factors.

The missing estimate for $Q_3$ illustrates the obstacle to carrying
over the automorphic part of the Lean proof. Detecting $v_2(c)=2$
in a representation $p=a^2+c^2$ uses splitting in a nonabelian
extension of $\mathbb Q$. Indeed the class group at conductor eight
is cyclic of order four: after choosing the odd coordinate as the real
coordinate, the slope $c/a\bmod8$ is an even residue, and doubling slope
2 under Gaussian multiplication gives slope 4 modulo 8. Complex
conjugation inverts this cyclic group. The resulting degree-eight
dihedral group distinguishes this test from restriction to
Dirichlet residue classes. No estimate for this new weight is provided
by the source's theorem about divisibility of its particular theta jets.

\section{The size of a negative excursion in the twisted prime sum}
\label{sec:osc}
The file \texttt{New/Reduction/Zeros/PrimeError.lean} computes
the poles of a smoothed prime sum used in the proof. We ask how far
its normalized error can go below zero when the auxiliary prime varies.
The answer for modulus one has an absolute uncertainty of less than
$0.047$ on either side of an explicit center. It also gives a concrete
question about relations between zeros on the critical line.

Fix a prime $d\equiv7\pmod 8$, let $\chi(n)=(n/d)$, and let
$\gamma_E$ be Euler's constant. With this notation the repository's
\texttt{primeRemainder} at modulus one is
\[
 R_d(x)=\sum_{n\le x}\Lambda(n)(1-\chi(n))(1/n-1/x)
       -\log x+1+\gamma_E-L'(1,\chi)/L(1,\chi).
\]
Put $B_0=2+\gamma_E-\log(4\pi)=0.0461914179\ldots$ and
$B_d=\log(d/\pi)-\gamma_E+2 L'(1,\chi)/L(1,\chi)$.

\paragraph{Excursion theorem.}
Under RH and Dirichlet GRH, for every such $d$,
\[
 \boxed{\quad -B_d-B_0\ \le\ \liminf_{x\to\infty}\sqrt x\,R_d(x)
                   \ \le\ -B_d+B_0.\quad}
\]
Moreover a single sequence $x_j\to\infty$ realizes
$\sqrt{x_j}R_d(x_j)\to B_0-B_d$ simultaneously for every $d$.
In particular the lower limit divided by $\log d$ tends to $-1$ as
$d\to\infty$ through these primes. The two limits here are taken in order,
first in $x$ for fixed $d$. The theorem makes no prediction about the
first $x$ at which a given excursion occurs.

\paragraph{The explicit formula.}
For $F=\zeta$ or $F=L(\cdot,\chi)$ write
\[
 H_F(u)=\sum_{\rho=1/2+i\gamma}
          \frac{e^{i\gamma u}}{1/4+\gamma^2}.
\]
Take all nontrivial zeros with multiplicities, including central zeros.
The zero count $O_d(T\log(2T))$ gives absolute uniform convergence.
Complex conjugation makes each series real. For fixed $d$,
\[
 \sqrt x\,R_d(x)=H_\zeta(\log x)-H_{L(\cdot,\chi)}(\log x)
                         +O_d(x^{-1/2}).                         \tag{O1}
\]
Indeed integrate $x^s A(s)/(s(s-1))$ on $\Re s=2$, with
$A=-\zeta'/\zeta+L'/L(\cdot,\chi)$ and measure $ds/(2\pi i)$.
Mellin inversion gives $\sum_{n\le x}\Lambda(n)(1-\chi(n))(x/n-1)$,
with zero weight at $n=x$. At $s=1$ the residue is
$x(\log x-1-\gamma_E+L'(1,\chi)/L(1,\chi))$. At $\rho$ it is
$x^\rho(m_\zeta(\rho)-m_L(\rho))/(\rho(1-\rho))$, including
coincident zeros. Shift the contour to $\Re s=-1/2$. The kernel
has a simple pole at zero since $\zeta(0)L(0,\chi)\ne0$ for this odd
primitive character. Its residue is a constant depending on $d$.
On the new vertical line the logarithmic derivatives are
$O_d(\log(|\Im s|+2))$ by the functional equations.
Horizontal edges may be sent to infinity along heights separated by
$\gg_d 1/\log T$ from both zero sets in $[T-1,T+2]$, using the
standard local zero count $O_d(\log T)$ and the local partial fraction
formula $F'/F(s)=\sum_{|\Im\rho-\Im s|\le1}(s-\rho)^{-1}
+O_d(\log T)$ in this fixed strip at large heights. Thus the edge
integrals vanish and the new integral is $O_d(x^{-1/2})$. Dividing
the shifted identity by $\sqrt x$ proves (O1).

\paragraph{Mass and recurrence.}
For the completed entire functions $\Xi_F$ (including the factor
$s(s-1)/2$ for zeta), the Hadamard product gives
\[
 \frac{\Xi'_F}{\Xi_F}(1)-\frac{\Xi'_F}{\Xi_F}(0)
    =\sum_{\rho}\frac{1}{\rho(1-\rho)}
    =H_F(0).
\]
The functional equation makes the left side twice the derivative
ratio at $1$. Substitution of the gamma factors gives $B_0$ for zeta
and $B_d$ for the odd quadratic character.
Consequently $|H_F(u)|\le H_F(0)$. This proves the lower bound in the box.
For finitely many real ordinates, simultaneous approximation returns
all their phases arbitrarily close to $1$ at arbitrarily large positive
times, regardless of rational relations. This follows by compactness:
the integer iterates in the finite torus have a convergent subsequence,
and differences of increasingly separated indices return to the identity.
Enumerate all ordinates for all the countably many characters and zeta,
then take $u_j>j$ approximating the first $j$ phases to within $1/j$.
Absolute convergence in each fixed series gives (O1) along $x_j=e^{u_j}$,
which proves the upper bound and the simultaneous assertion. There is
even a single sequence of integer exponents $n_j\to\infty$ such that
$\sqrt{a^{n_j}}R_d(a^{n_j})\to B_0-B_d$ for every integer $a\ge2$
and every eligible $d$: apply the same integer-iterate argument to the
countably many frequencies $\gamma\log a$. In particular the boxed bounds
hold along powers of any given integer base. Finally
the GRH estimate $L'(1,\chi)/L(1,\chi)=O(\log\log d)$ gives the conductor
asymptotic. Ihara, Murty and Shimura prove this estimate and an explicit
formula for this same ramp kernel in
\href{https://www.impan.pl/shop/en/publication/transaction/download/product/82940}
{\emph{On the logarithmic derivatives of Dirichlet $L$-functions at $s=1$}},
Acta Arith.\ 137 (2009), Corollary 3.3.2 and \S4.2. The contribution
here is the two-sided localization and common recurrence for the
difference used by the Lean development.

\paragraph{How much information is lost in the interval?}
The lower endpoint is attained as a limiting value precisely when every
finite integer relation among the positive ordinates of zeta and of
$L(\cdot,\chi)$ has even total coefficient on the zeta ordinates.
Treat each occurrence as labeled by its function, so a shared positive
ordinate already violates the condition. To see the assertion, equality
forces each positive-weight cosine of zeta to approach $-1$ and each
cosine of $L$ to approach $+1$: their individually nonnegative deficits
sum to zero in the limit. Zeta has no central zero, since the alternating
eta series is positive at $1/2$. Central zeros of $L$ already contribute
with the required sign. On each finite set of frequencies, Kronecker's
approximation theorem identifies the closure of the real flow with the
subtorus satisfying its integer relations. The desired sign vector
belongs to that closure exactly under the stated parity condition.
A diagonal sequence and the summable tails give sufficiency for all
ordinates. This identifies the additional arithmetic input that would
give an exact lower endpoint. Whether the parity condition holds for
these spectra is unresolved here.

\paragraph{Numerical examples.}
For $d=7$ the endpoints are approximately $-0.30142740$ and
$-0.20904456$. For $d=23$ they are $-0.61549819$ and $-0.52311535$.
These values evaluate the proven bounds; they do not estimate a lower
limit by finite sampling. The companion script uses the finite
Hurwitz-zeta formula at zero and the functional equation. In detail, put
$A=-d^{-1}\sum_{a=1}^{d-1}a\chi(a)=L(0,\chi)$. Then
\[
 \frac{L'(1,\chi)}{L(1,\chi)}
 =\log(2\pi)+\gamma_E
  -\frac{1}{A}\sum_{a=1}^{d-1}\chi(a)\log\Gamma(a/d).
\]


\paragraph{Sources and comparison.} For ideal-class characters and
their Hecke functions see Duke, Friedlander and Iwaniec,
\href{https://www.math.ucla.edu/~wdduke/preprints/classgroup.pdf}{\emph{Class
group $L$-functions}}. For the form--field correspondence see Cox,
\href{https://www.math.utoronto.ca/~ila/Cox-Primes_of_the_form_x2%2Bny2.pdf}{\emph{Primes
of the form $x^2+ny^2$}}, especially \S9. Suzuki gives related
\href{https://doi.org/10.1007/s11139-025-01238-9}{one-sided criteria}
for zeta and Dirichlet functions. The references and pinned source
links identify which inputs are imported. The exact Fourier
optimization and the equivalences in this note are proved in the text.
\end{document}
